\documentclass[11pt,a4paper]{article}
\usepackage[italian]{babel}
\usepackage[margin=2.5cm]{geometry}
\usepackage{parskip}
\usepackage{amsmath, amssymb, amsthm}
\usepackage{booktabs}
\usepackage{array}
\usepackage{xcolor}
\usepackage{needspace}
\usepackage{enumitem}
\usepackage{listings}
\usepackage{tikz}
\usepackage[most]{tcolorbox}
\usepackage[hidelinks]{hyperref}
\usetikzlibrary{decorations.pathreplacing, shapes.geometric, arrows.meta, positioning}

% Niente vedove/orfane: i paragrafi non lasciano righe isolate a fine/inizio pagina
\widowpenalty=10000
\clubpenalty=10000

% Elenchi più compatti
\setlist{itemsep=2pt, parsep=0pt, topsep=4pt, partopsep=0pt}

% --- Riquadri con bordo nero, senza numero (si spezzano tra due pagine) ---
% Uso: \begin{definizione} ... \end{definizione}
%      \begin{teorema}[Nome] ... \end{teorema}  (il nome è facoltativo)
\tcbset{nota/.style={enhanced jigsaw, breakable, colback=white,
  colframe=black, boxrule=0.5pt, sharp corners}}
\NewTColorBox{definizione}{}{nota,
  before upper={\textbf{Definizione.}\ }}
\NewTColorBox{teorema}{o}{nota,
  before upper={\textbf{Teorema\IfValueT{#1}{ (#1)}.}\ }}
\NewTColorBox{proposizione}{o}{nota,
  before upper={\textbf{Proposizione\IfValueT{#1}{ (#1)}.}\ }}
\NewTColorBox{proprieta}{o}{nota,
  before upper={\textbf{Proprietà\IfValueT{#1}{ (#1)}.}\ }}
\NewTColorBox{assioma}{o}{nota, boxrule=1pt,
  before upper={\textbf{Assioma\IfValueT{#1}{ (#1)}.}\ }}

% --- Ambienti semplici (senza riquadro) ---
% Il titolo sta in cima al blocco, anche se il contenuto inizia con un riquadro o
% con codice; e non resta mai da solo in fondo a una pagina.
% Uso: \begin{esempio}[Titolo facoltativo] ... \end{esempio}
\newcommand{\newrunin}[2]{%
  \NewDocumentEnvironment{#1}{o}{%
    \par\Needspace{4\baselineskip}\addvspace{\medskipamount}\noindent
    \textbf{#2}\IfValueT{##1}{ (##1)}\textbf{.}\ \ignorespaces}%
  {\par\addvspace{\medskipamount}}}
\newrunin{esempio}{Esempio}
\newrunin{esercizio}{Esercizio}
\NewTColorBox{osservazione}{}{nota,
  before upper={\textbf{Osservazione.}\ }}

% V e F nelle tavole di verità
\newcommand{\V}{\textsf{V}}
\newcommand{\F}{\textsf{F}}

% --- Codice C ---
% Uso: \begin{codice} ... \end{codice}      codice C numerato
%      \begin{comando} ... \end{comando}    comandi da terminale
%      \begin{risultato} ... \end{risultato} output di un programma
%      \lstinline|printf("ciao");|           codice dentro al testo
\lstdefinestyle{c}{
  language=C,
  basicstyle=\ttfamily\small,
  keywordstyle=\color{blue}\bfseries,
  commentstyle=\color{gray}\itshape,
  stringstyle=\color{red!70!black},
  numbers=left,
  numberstyle=\tiny\color{gray},
  numbersep=8pt,
  columns=flexible,
  keepspaces=true,
  breaklines=true,
  showstringspaces=false,
  tabsize=4,
  morekeywords={include, define},
  % lettere accentate nei commenti
  literate={à}{{\`a}}1 {è}{{\`e}}1 {é}{{\'e}}1 {ì}{{\`i}}1
           {ò}{{\`o}}1 {ù}{{\`u}}1 {È}{{\`E}}1 {À}{{\`A}}1
}
\lstset{style=c}
\newtcblisting{codice}{listing only, nota, left=6mm,
  listing options={style=c}}
\newtcblisting{comando}{listing only, nota,
  listing options={basicstyle=\ttfamily\small, numbers=none, language={},
    columns=flexible, keepspaces=true, breaklines=true}}
\newtcblisting{risultato}{listing only, enhanced jigsaw, breakable,
  colback=black!5, colframe=black, boxrule=0.5pt, sharp corners,
  listing options={basicstyle=\ttfamily\small, numbers=none, language={},
    columns=flexible, keepspaces=true, breaklines=true}}

% --- Diagrammi di flusso (simboli standard) ---
\tikzset{
  terminale/.style = {ellipse, draw, minimum width=2.5cm, minimum height=0.9cm},
  io/.style        = {trapezium, trapezium left angle=70, trapezium right angle=110,
                      draw, minimum width=2.5cm, minimum height=0.9cm},
  processo/.style  = {rectangle, draw, minimum width=2.5cm, minimum height=0.9cm},
  decisione/.style = {diamond, draw, aspect=2, minimum width=2.5cm},
  freccia/.style   = {thick, -{Stealth}}
}

\title{Selezione, salti e typedef}
\author{Marco Cerbella}
\date{Lezione 14 -- 26 ottobre 2026}

\begin{document}
\maketitle

\section{Istruzioni di selezione}

Un'istruzione di selezione può dirigere il flusso di esecuzione del programma lungo percorsi diversi a seconda di una condizione data. In C ci sono due istruzioni di selezione:
\begin{enumerate}
    \item \lstinline|if|;
    \item \lstinline|switch|.
\end{enumerate}

\subsection{Istruzione if}

Un'istruzione \lstinline|if| ha la forma
\[
\texttt{if ( espressione ) istruzione1 [ else istruzione2 ]}
\]
La clausola \lstinline|else| è facoltativa. L'espressione viene valutata per prima, per determinare quale delle due istruzioni viene eseguita; deve avere tipo scalare. Se il suo valore è vero (cioè diverso da $0$) viene eseguita \texttt{istruzione1}; altrimenti viene eseguita \texttt{istruzione2}, se presente.

\begin{codice}
if (testExpression)
{
    // corpo dell'if
}
// istruzione subito dopo l'if
\end{codice}
\begin{codice}
if (testExpression)
{
    // corpo dell'if
}
else
{
    // corpo dell'else
}
// istruzione subito dopo l'if...else
\end{codice}

\subsubsection*{If annidati}

Se più istruzioni \lstinline|if| sono annidate, una clausola \lstinline|else| appartiene sempre all'ultimo \lstinline|if| (allo stesso livello di annidamento dei blocchi) che non ha ancora un \lstinline|else|:
\begin{codice}
if (n > 0)
    if (n % 2 == 0)
        puts("n is positive and even");
    else    // alternativa all'ultimo if
        puts("n is positive and odd");
\end{codice}
Una clausola \lstinline|else| può essere assegnata a un \lstinline|if| diverso racchiudendo in un blocco l'ultimo \lstinline|if| che non deve avere un \lstinline|else|:
\begin{codice}
if (n > 0)
{
    if (n % 2 == 0)
        puts("n is positive and even");
}
else        // alternativa al primo if
    puts("n is negative or zero");
\end{codice}

\subsubsection*{Catena else if}

Per scegliere una tra più di due istruzioni alternative si possono mettere a cascata degli \lstinline|if| in una catena \lstinline|else if|. Ogni nuovo \lstinline|if| è semplicemente annidato nella clausola \lstinline|else| di quello precedente:
\begin{codice}
double spec = 10.0, measured = 10.7, diff;
diff = measured - spec;
if (diff >= 0.0 && diff < 0.5)
    printf("Rounding down: %.2f\n", diff);
else if (diff >= 0.5 && diff < 1)
    printf("Rounding up: %.2f\n", diff);
else
    printf("Difference greater than 1!\n");
\end{codice}
Appena una di queste espressioni è vera, viene eseguita l'istruzione corrispondente e le altre vengono scartate. Se nessuna delle condizioni è vera, viene eseguita la clausola \lstinline|else| dell'ultimo \lstinline|if|, se presente.

\begin{osservazione}
Nelle slide la seconda \lstinline|printf| stampava ancora «Rounding down»: qui è stata corretta in «Rounding up», coerentemente con l'intervallo $[0{,}5,\,1)$.
\end{osservazione}

\subsection{Istruzione switch}

Un'istruzione \lstinline|switch| fa saltare il flusso di esecuzione a una tra più istruzioni, a seconda del valore di un'espressione intera:
\[
\texttt{switch ( espressione ) istruzione}
\]
\texttt{espressione} ha tipo intero e \texttt{istruzione} è il \emph{corpo dello switch}, che contiene etichette \lstinline|case| e al più un'etichetta \lstinline|default|. L'espressione viene valutata una sola volta e confrontata con le espressioni costanti nelle etichette \lstinline|case|. Se il valore coincide con una delle costanti, il flusso del programma salta all'istruzione che segue quell'etichetta. Se nessuna costante coincide, il programma continua all'etichetta \lstinline|default|, se c'è.

\begin{codice}
switch (menu())   // salta in base al risultato di menu()
{
    case 'A': action1();     // esegue l'azione 1
              break;         // non esegue nessun'altra "azione"
    case 'B': action2();     // esegue l'azione 2
              break;         // non esegue l'azione di default
    default:  putchar('\a'); // nessun comando riconosciuto
}
\end{codice}

\subsubsection*{Intervalli}

Per trattare un intervallo di valori con la stessa azione, si possono scrivere più etichette consecutive:
\begin{codice}
switch (value)
{
    case 1:
    case 2:
    case 3:
    case 4:
    case 5:
    case 6:
    case 7:
    case 8:
        printf("Hello, 1 to 8\n");
        break;
    default:
        printf("Hello, default\n");
        break;
}
\end{codice}
Esiste anche un'estensione di GCC, non standard, che permette di scrivere \lstinline|case 1 ... 8:|. Un'alternativa è mettere nel \lstinline|default| un \lstinline|if|:
\begin{codice}
switch (value) {
    case 1:
        /* do 1 */
        break;
    case 2:
        /* do 2 */
        break;
    default:
        if (value >= 3 && value <= 8)
            printf("Hello, 3 to 8\n");
}
\end{codice}
Se però si deve lavorare con molti intervalli, conviene usare gli \lstinline|if|.

\subsubsection*{Senza break: fall-through}

Senza \lstinline|break| l'esecuzione prosegue nelle etichette successive.
\begin{esempio}
\begin{codice}
#include <stdio.h>
int main()
{
    int Grade = 'A';
    switch (Grade)
    {
        case 'A': printf("Excellent\n");
        case 'B': printf("Good\n");
        case 'C': printf("OK\n");
        case 'D': printf("Mmmmm....\n");
        case 'F': printf("You must do better than this\n");
        default : printf("What is your grade anyway?\n");
    }
}
\end{codice}
\begin{risultato}
Excellent
Good
OK
Mmmmm....
You must do better than this
What is your grade anyway?
\end{risultato}
\end{esempio}

\section{Salti incondizionati}

Le istruzioni di salto interrompono l'esecuzione sequenziale delle istruzioni, in modo che l'esecuzione continui in un altro punto del programma. Un salto distrugge le variabili automatiche se la destinazione del salto è fuori dal loro scope. In C ci sono quattro istruzioni che causano salti incondizionati: \lstinline|break|, \lstinline|continue|, \lstinline|goto| e \lstinline|return|.

\subsection{break}

L'istruzione \lstinline|break;| può comparire solo nel corpo di un ciclo o di uno \lstinline|switch| e causa un salto alla prima istruzione dopo il ciclo o lo \lstinline|switch| che la contiene immediatamente. Si può quindi usare \lstinline|break| per terminare l'esecuzione di un ciclo in qualunque punto del corpo.

\begin{esempio}
\begin{codice}
#include <stdio.h>
int main()
{
    int i;
    double number, sum = 0.0;

    for (i = 1; i <= 10; ++i)
    {
        printf("Enter a n%d: ", i);
        scanf("%lf", &number);

        // se l'utente inserisce un numero negativo, il ciclo termina
        if (number < 0.0)
        {
            break;
        }
        sum += number;
    }
    printf("Sum = %.2lf", sum);
    return 0;
}
\end{codice}
Con input $2{,}4$, $4{,}5$, $3{,}4$, $-3$ l'output è \texttt{Sum = 10.30}.
\end{esempio}

\subsection{continue}

L'istruzione \lstinline|continue;| si può usare solo nel corpo di un ciclo e fa saltare il resto dell'iterazione corrente.
\begin{itemize}
    \item In un ciclo \lstinline|while| o \lstinline|do ... while| il programma salta alla successiva valutazione dell'espressione di controllo.
    \item In un ciclo \lstinline|for| il programma salta alla successiva valutazione della terza espressione dell'istruzione \lstinline|for|, che contiene le operazioni eseguite dopo ogni iterazione.
\end{itemize}

\begin{esempio}
Stesso programma di prima, ma con \lstinline|continue| al posto di \lstinline|break|:
\begin{codice}
for (i = 1; i <= 10; ++i)
{
    printf("Enter a n%d: ", i);
    scanf("%lf", &number);

    // se l'utente inserisce un numero negativo, si passa all'iterazione successiva
    if (number < 0.0)
    {
        continue;
    }
    sum += number;
}
\end{codice}
Con input $1{,}1$, $2{,}2$, $5{,}5$, $4{,}4$, $-3{,}4$, $-45{,}5$, $34{,}5$, $-4{,}2$, $-1000$, $12$ i negativi vengono saltati e l'output è \texttt{Sum = 59.70}.
\end{esempio}

\subsection{goto}

L'istruzione \lstinline|goto| causa un salto incondizionato a un'altra istruzione nella stessa funzione. La destinazione del salto è specificata dal nome di un'\emph{etichetta}:
\[
\texttt{goto nome\_etichetta;}
\]
Un'etichetta è un nome seguito da due punti: \texttt{nome\_etichetta: istruzione}. Le etichette hanno un proprio spazio dei nomi, quindi possono avere gli stessi nomi di variabili o tipi senza causare conflitti. Le etichette possono essere messe prima di qualunque istruzione, e un'istruzione può avere più etichette. Servono solo come destinazioni dei \lstinline|goto| e non hanno alcun effetto se l'istruzione etichettata viene raggiunta nel normale corso dell'esecuzione sequenziale.

\begin{codice}
void calculate(int a)
{
    if (a < 1 || a > 5)
        goto here;
    printf("a between 1 and 5!\n");
    return;
here: printf("a < 1 or a > 5!\n");
    return;
}
\end{codice}

\begin{osservazione}
Poiché il codice che fa largo uso di \lstinline|goto| è difficile da leggere, conviene usarli solo quando offrono un vantaggio evidente, come l'uscita rapida da cicli profondamente annidati. Qualunque programma C che usa \lstinline|goto| si può scrivere anche senza.
\end{osservazione}

Lo \emph{spaghetti code} è un'espressione spregiativa per il codice sorgente con una struttura di controllo complessa e ingarbugliata, soprattutto quando usa molti \lstinline|goto|, eccezioni, thread o altri costrutti di salto «non strutturati» (tipico di assembly, Fortran, BASIC).

\subsection{return}

Si veda la lezione sulle funzioni: \lstinline|return| è anch'esso un'istruzione di salto incondizionato, come \lstinline|goto|.

\section{typedef}

Il modo più semplice di usare tipi con nomi complessi è dichiarare dei \emph{sinonimi} semplici per essi. In una dichiarazione che inizia con la parola chiave \lstinline|typedef|, ogni dichiaratore definisce un identificatore come sinonimo del tipo specificato. L'identificatore è poi detto \emph{typedef name} per quel tipo. A parte la parola chiave \lstinline|typedef|, la sintassi è esattamente la stessa di una dichiarazione di un oggetto o di una funzione del tipo specificato.

\begin{esempio}
\begin{codice}
typedef unsigned int UINT;
typedef enum state {DEAD, ALIVE} State;

UINT ui = 10;
State s = DEAD;
\end{codice}
Nello scope di queste dichiarazioni \lstinline|UINT| è sinonimo di \lstinline|unsigned int| e \lstinline|State| è sinonimo di \lstinline|enum state|. La variabile \lstinline|ui| ha tipo \lstinline|unsigned int| e \lstinline|s| ha tipo \lstinline|enum state|.
\end{esempio}

\section{Riferimenti}

Sul libro: sezione 10.6 (\lstinline|typedef|).

\end{document}
